\documentclass[a4paper, 10pt]{article}
\usepackage{scrextend}
\changefontsizes{9pt}
\input{../course-config}
\usepackage{../vendor/template/CEDT-Assignment-style}

\begin{document}
\subject
\asgntitle{X}{}{Week x - x}{\StudentID\ \StudentName}{\DefaultCollaborators}

% This is a lean starting point -- see
% vendor/template/templates/style-guide/style-guide.pdf (source:
% .../style-guide.tex) for a worked example of every macro/environment
% CEDT-Assignment-style.sty provides (theorems, math shortcuts, tables, figures,
% multi-language codingbox, appendix, ...). Nothing here points at a file that
% won't exist in a freshly scaffolded assignment folder -- keep it that way, or
% move file-dependent examples (\fig, \myappendix, ...) to your own solved code.

% ================================================================================ %
\section{PUT THE SECTION HERE}
% ================================================================================ %

% ================================================================================ %
%                                    Problem 0x                                    %
% ================================================================================ %
\begin{problem}
Here's problem statement.
\begin{codingbox}
    print("You can insert code like this")
\end{codingbox}
You can insert sub-problem
\begin{subproblems}
	\item Here's sub-problem a.
	\item Here's sub-problem b.
	\item Here's sub-problem c.
\end{subproblems}
\end{problem}

\begin{solution}
	The solution goes here.
\end{solution}
% ================================================================================ %

% ================================================================================ %
%                                    Problem 0x                                    %
% ================================================================================ %
\begin{problem}
Here's the problem statement for \textbf{MANDATORY} problem.
\end{problem}

\begin{tosubmit}
	Here's the solution for the \textbf{MANDATORY} problem. \\
	ABC.
\end{tosubmit}
% ================================================================================ %

% ================================================================================ %
%                                    Problem 0x                                    %
% ================================================================================ %
\begin{problem}
Here's the problem statement for \textbf{PROVING} problem.
\end{problem}

\begin{proofbox}
	Here's the proof for the \textbf{PROVING} problem.
\end{proofbox}
% ================================================================================ %

% ================================================================================ %
%                                    Problem 0x                                    %
% ================================================================================ %
\begin{problem}
Here's a \textbf{multiple choice} problem statement. Use \verb|\choices| to lay
out the options; it defaults to 4 choices and picks 4, 2, or 1 columns on its
own depending on how long the choices are.
\begin{subproblems}
	\item Short choices sit four to a row.
	      \choices{True, False, Maybe, Unknown}
\end{subproblems}
\end{problem}
% ================================================================================ %

% ================================================================================ %
%                                 Writing in Thai                                  %
% ================================================================================ %
% This template builds with XeLaTeX (the .latexmkrc here sets that up). Thai text
% must be wrapped: \textthai{...} for a few words inline, the thai environment
% for whole paragraphs -- works anywhere (problem statements, \choices, section
% titles, \asgntitle). See the style guide for the full reference.
\begin{problem}
\textthai{ตัวอย่างโจทย์ภาษาไทย} -- Thai and English mix freely.
\end{problem}

\begin{solution}
	\begin{thai}
		เขียนคำตอบเป็นภาษาไทยทั้งย่อหน้าได้ ตัวตัดคำจัดการขึ้นบรรทัดใหม่ให้เอง
		และแทรกคณิตศาสตร์อย่าง $a^2 + b^2 = c^2$ ได้ตามปกติ
	\end{thai}
\end{solution}
% ================================================================================ %

\end{document}
